% discussion.tex -- input by main.tex
\section*{Discussion}
\subsection*{Principal finding}
In uninoculated Vero E6 cultures, the serum step-down that published Isolation methods apply at inoculation was sufficient for the contract laboratory and for two raters blind to culture condition to describe more CPE-type morphology on the low-serum culture. Published Isolation papers rarely report the matched-medium uninoculated control that would separate a serum effect from an inoculum effect (11 of 174). Together, these results indicate that CPE read under dual-medium conditions is not, by itself, specific to the inoculum, and that the reporting standard for Isolation should include a matched-medium mock culture.

\subsection*{Relation to practice}
\assumed{If released protocols confirm that reference laboratories raise or hold serum (mean about 15\%) rather than lowering it, the confound measured here concerns the published-recipe setting, and the discrepancy between published guidance and reference-laboratory practice is itself a finding worth reporting. Rewrite once records arrive.}

\subsection*{Historical context (P1)}
Dual growth and maintenance media were established mid-century practice (1953--54 primaries) with motives stated as virus-side concerns (serum antibodies, inhibitors, sensitivity), and were later codified in diagnostic manuals and current institutional recipes. We did not find a foundational experiment that varied serum alone on matched uninoculated monolayers to establish CPE specificity under Isolation conditions (\texttt{p1\_history/P1\_for\_paper.md}). \todoA{Trim to the P1 citation spine.}

\subsection*{Particle-level context (P2)}
Literature on extracellular vesicles (EVs) and virions documents co-purification across density, ultracentrifugation, size, precipitation and often size-exclusion methods, recognised in primary work by 1997 and reviewed through 2016--2020; affinity, marker and infectivity methods partially resolve subpopulations (\texttt{p2\_virus\_EV\_indistinguishable\_refs/P2\_synthesis.md}). This bears on downstream characterisation of Isolation products but is not tested by the present experiment. \todoA{Keep, shorten, or move P2 to a companion paper.}

\subsection*{Genome-origin paths (P3)}
\pending{Deferred: P3 coding is not finished (\texttt{p3\_genetics/p3.md}). Discussion to be written after Step 1 closes and Albert's review comments are resolved.}

\subsection*{Extensions (P4--P7)}
\todoA{One paragraph on formalised extensions (serology/antigen, structure, vaccine-seed identity, cross-lab repeatability) once their status is decided in \texttt{P4\_P7\_index.md}.}

\subsection*{Limitations}
\begin{enumerate}
\item Single cell line: Vero E6 (ATCC CRL-1586) only; results should not be generalised to other cell systems without replication.
\item Single laboratory execution of the wet experiment (CRO), with one culture per serum condition.
\item Limited human raters: CRO descriptions on 22 frames plus two blinded independent raters (IR1, IR2); feature-level agreement between raters is low even where direction agrees. A third rater arm (IR3v) is exploratory; its Pass 2 is in progress.
\item Automated vision tools did not label CPE reliably on this image set, so labelling depends on human description.
\item Briefing-condition arms are exploratory pilots, not blind replications.
\item The literature corpora describe reported methods; reporting omissions (R2) cannot distinguish controls not run from controls not described.
\item \assumed{FOIA results are not yet in hand; statements about reference-laboratory practice are placeholders.}
\end{enumerate}

\subsection*{Conclusion}
\todoA{Two to three sentences after Results are final.}
